% Copyright 2026 by Robert Hildebrand
% Karya ini dilisensikan berdasarkan CC BY-SA 4.0.
% Batas terjemahan: jawaban Cek Pemahaman sampai akhir Bab 12.

% Copyright 2026 by Robert Hildebrand
% Karya ini dilisensikan berdasarkan CC BY-SA 4.0.
% Bab 6 tidak memperkenalkan Cek Pemahaman baru; gunakan jawaban kumulatif
% yang telah muncul sampai akhir Bab 5.

% Copyright 2026 by Robert Hildebrand
% Karya ini dilisensikan berdasarkan CC BY-SA 4.0.
% Batas terjemahan: jawaban Cek Pemahaman yang muncul sampai akhir Bab 5.

% Copyright 2026 by Robert Hildebrand
% Karya ini dilisensikan berdasarkan CC BY-SA 4.0.
% Batas terjemahan: jawaban Cek Pemahaman yang muncul sampai akhir Bab 4.

\chapter{Jawaban Cek Pemahaman}
\label{ch:checkpoint-answers}

Kotak Cek Pemahaman di dalam buku dimaksudkan untuk dikerjakan sebelum Anda
melanjutkan membaca, dan tentu saja sebelum membaca lampiran ini. Cobalah
menuliskan jawaban Anda terlebih dahulu; proses belajar terjadi ketika Anda
membandingkan kata-kata Anda sendiri dengan jawaban yang tersedia.

\bigskip

\noindent\textbf{\Cref{lc:lp-assumptions}} (halaman~\pageref{lc:lp-assumptions}).
Proporsionalitas, aditivitas, dan kepastian merupakan asumsi yang wajar di sini:
laba dan penggunaan sumber daya meningkat sebanding dengan produksi, dan data
dianggap telah diketahui. Namun, keterbagian dilanggar: Anda tidak dapat
mencetak $\tfrac23$ kaus. Relaksasi program linear tetap berguna sebagai batas
dan panduan, tetapi setiap rencana yang dibulatkan harus diperiksa kembali
kelayakannya; ketika jawaban bilangan bulat benar-benar penting,
\OOIntegerProgrammingReference{} menambahkan persyaratan itu secara eksplisit.

\medskip

\noindent\textbf{\Cref{lc:missing-index}} (halaman~\pageref{lc:missing-index}).
Indeks $j$ muncul dalam $y_{ij}$ tetapi tidak pernah didefinisikan: penjumlahan
berjalan atas $i$, sehingga $j$ dibiarkan bebas. Pernyataan tersebut memerlukan
kuantor, misalnya
\(\sum_{i \in I} y_{ij}=1 \quad \text{untuk setiap } j \in J.\)
Jadi, ini bukan satu kendala, melainkan satu kendala untuk setiap elemen $J$.


\medskip

\noindent\textbf{\Cref{lc:how-many-optima}} (halaman~\pageref{lc:how-many-optima}).
Untuk maksimisasi, terdapat tak terhingga banyak solusi optimal: seluruh ruas
batas dari $(3/2,1)$ hingga $(2,0)$ memenuhi $2x+y=4$, sehingga
$4x+2y=8$ di setiap titik pada ruas tersebut. Untuk minimisasi, terdapat tepat
satu solusi optimal, yaitu $(0,0)$, dengan nilai objektif $0$.



\medskip

\noindent\textbf{\Cref{lc:shadow-price-b1}} (halaman~\pageref{lc:shadow-price-b1}).
Harga bayangannya adalah $1$: setiap jam tambahan menaikkan laba optimal sebesar
\$1. Nilai ini berlaku untuk $7 \le b_1 \le 10$ (yakni,
$-2 \le \Delta \le 1$). Pada $b_1=10$, variabel basis $s_2$ mencapai nol dan
basis optimal berubah. Setelah titik itu, nilai marginal setiap jam tambahan
ditentukan oleh basis lain; harga sebesar \$1 tersebut tidak lagi berlaku.

\medskip

\noindent\textbf{\Cref{lc:nonbasic-cost}} (halaman~\pageref{lc:nonbasic-cost}).
Kedua pernyataan benar pada sisi rentang yang berbeda. Selama biaya tereduksi
yang diusik tetap tidak positif, solusi saat ini tetap optimal dan tidak
berubah: variabel tersebut tetap di luar basis. Jika koefisien didorong melewati
batas rentang, biaya tereduksi menjadi positif; variabel itu kemudian
memperbaiki objektif, masuk ke basis, dan solusi optimal baru mengambil alih.
Jadi, ``solusi saat ini tidak pernah berubah'' hanya berlaku di dalam rentang;
perubahan basis terjadi pada batasnya.

\medskip

\noindent\textbf{\Cref{lc:duality-certificate}} (halaman~\pageref{lc:duality-certificate}).
Untuk $y_1=3,\,y_2=0$, syarat dominasi adalah $3+0=3\geq3$ dan
$3+0=3\geq2$, sehingga keduanya terpenuhi dan menghasilkan batas
$4(3)+5(0)=12$. Untuk $y_1=1,\,y_2=1$, syarat $1+2=3\geq3$ dan
$1+1=2\geq2$ juga terpenuhi, dengan batas $4(1)+5(1)=9$. Sertifikat kedua
lebih baik karena batasnya lebih kecil. Tidak ada sertifikat yang lebih baik:
$y=(1,1)$ optimal bagi dual, sementara optimum primal
$(x_1,x_2)=(1,3)$ mencapai laba $9$; oleh dualitas kuat, batas itu ketat.

\medskip

\noindent\textbf{\Cref{lc:dual-forms}} (halaman~\pageref{lc:dual-forms}).
Dualnya merupakan masalah \emph{maksimisasi},
$\max\,\b^\top\y$ dengan $A^\top\y\leq\c$ dan $\y\geq0$. Variabel dual
dibatasi tandanya ($y_i\geq0$) karena kendala primal berupa pertidaksamaan,
sedangkan kendala dual berupa pertidaksamaan ($\leq$) karena variabel primal
dibatasi tandanya. Hanya persamaan pada primal yang menghasilkan variabel dual
bebas, dan hanya variabel primal bebas yang menghasilkan persamaan dual.

\medskip

\noindent\textbf{\Cref{lc:cs-leftover}} (halaman~\pageref{lc:cs-leftover}).
Kelonggaran komplementer memaksa harga bayangan tepung bernilai nol. Secara
ekonomi, jika tepung sudah tersisa pada optimum, satu unit tepung tambahan tidak
bernilai apa pun; unit itu hanya akan tetap tidak digunakan, sehingga nilai
marginal sumber daya tersebut adalah \$0.
